\documentclass[11pt]{article}

\usepackage[margin=1in]{geometry}
\usepackage{amsmath, amssymb, amsthm, bm, mathtools}
\usepackage{booktabs}
\usepackage{enumitem}
\usepackage{hyperref}
\usepackage{xcolor}

\newcommand{\E}{\mathbb{E}}
\newcommand{\Var}{\operatorname{Var}}
\newcommand{\Cov}{\operatorname{Cov}}
\newcommand{\tr}{\operatorname{tr}}
\newcommand{\R}{\mathbb{R}}
\newcommand{\dd}{\mathrm{d}}
\newcommand{\mc}[1]{\mathcal{#1}}
\newcommand{\mb}[1]{\mathbf{#1}}
\newcommand{\bs}[1]{\boldsymbol{#1}}

\newcommand{\xt}{x_t}
\newcommand{\pit}{\pi_t}
\newcommand{\ipol}{i_t}
\newcommand{\rn}{r^n_t}
\newcommand{\sigat}{\sigma_{a,t}}
\newcommand{\epsa}{\varepsilon_{a,t}}
\newcommand{\etasig}{\eta_{\sigma,t}}

\newcommand{\omegax}{\omega^x_t}
\newcommand{\omegapi}{\omega^\pi_t}
\newcommand{\omegai}{\omega^i_t}
\newcommand{\omegavec}{\bs{\omega}_t}

\newtheorem{remark}{Remark}
\newtheorem{definition}{Definition}


\title{A Minimal New Keynesian Accounting Note for TFP Uncertainty Wedges}
\author{Caught in the Crossfire Research Notes}
\date{\today}

\begin{document}

\maketitle

\begin{abstract}
This note develops a deliberately small New Keynesian accounting framework for
thinking about uncertainty shocks. The model has three reduced-form equilibrium
conditions: an IS equation, a New Keynesian Phillips curve, and a monetary policy
rule. The only primitive source of second-moment variation is stochastic
volatility in total factor productivity. The purpose is not to build the final
small-open-economy model. Instead, the purpose is to check whether a volatility
shock to productivity can be represented as a small set of wedges, in the spirit
of Business Cycle Accounting, after a second-order expansion of the underlying
nonlinear equilibrium conditions.
\end{abstract}

\tableofcontents

\clearpage

\section{Motivation and Scope}

The immediate goal is to create a controlled algebraic environment for thinking
about second-moment shocks before moving to the full small-open-economy
nonlinear BCA framework. The exercise starts from the smallest useful New
Keynesian system:
\begin{enumerate}[label=(\roman*)]
    \item an intertemporal demand equation,
    \item a price-setting equation, and
    \item a monetary policy rule.
\end{enumerate}
The only primitive uncertainty shock is stochastic volatility in productivity.
The question is whether the effects of that volatility shock can be organized as
wedges in the three equations.

This note should be read as an accounting device rather than a final structural
model. A first-order New Keynesian model is certainty equivalent, so a pure
volatility shock has no independent effect on allocations. A second-order
expansion of the underlying nonlinear equilibrium conditions produces risk
correction terms, and those terms can be collected as wedges. In a fully solved
DSGE model with stochastic volatility, third-order perturbation is usually
required for volatility shocks to enter the policy functions as independent
state variables. The second-order exercise here is therefore a diagnostic step:
it clarifies where the wedges would appear and what restrictions would be
implied by the assumption that TFP uncertainty is the only primitive source of
second-moment variation.

The intended bridge to the larger project is:
\begin{align*}
    \text{TFP volatility shock}
    &\Rightarrow
    \text{risk correction terms}
    \\
    &\Rightarrow
    \text{NK wedges}
    \\
    &\Rightarrow
    \text{BCA-style counterfactuals}.
\end{align*}

The note deliberately does not impose the later two-factor empirical structure
from the OI-SVMVAR. It focuses on a single source of uncertainty so that the
mapping from volatility to wedges is transparent.

\section{Baseline Three-Equation NK Model}

Let \(x_t\) denote the output gap, \(\pi_t\) inflation, \(i_t\) the nominal
policy rate, and \(r^n_t\) the natural real rate. The standard first-order
three-equation New Keynesian system is
\begin{align}
    x_t
    &=
    \E_t x_{t+1}
    -
    \frac{1}{\gamma}
    \left(
        i_t - \E_t \pi_{t+1} - r^n_t
    \right),
    \label{eq:baseline_is}
    \\
    \pi_t
    &=
    \beta \E_t \pi_{t+1}
    +
    \kappa x_t,
    \label{eq:baseline_nkpc}
    \\
    i_t
    &=
    \rho_i i_{t-1}
    +
    (1-\rho_i)
    \left(
        \phi_\pi \pi_t + \phi_x x_t
    \right)
    +
    \varepsilon^i_t.
    \label{eq:baseline_taylor}
\end{align}
The parameters satisfy \(\gamma>0\), \(0<\beta<1\), \(\kappa>0\), and
\(0\leq \rho_i<1\). In the baseline accounting exercise, monetary policy is not
a primitive uncertainty source, so \(\varepsilon^i_t\) is set to zero unless the
policy rule is used as a residual measurement equation.

Productivity enters the reduced-form NK system through the natural allocation.
For example, in a flexible-price economy with log productivity \(a_t\), natural
output and the natural real rate can be summarized as
\begin{align}
    y^n_t &= \chi_a a_t,
    \label{eq:natural_output}
    \\
    r^{n,1}_t
    &=
    \rho
    +
    \gamma \E_t \Delta y^n_{t+1},
    \label{eq:first_order_natural_rate}
\end{align}
where \(\rho=-\log\beta\), \(\chi_a\) is the elasticity of natural output with
respect to productivity, and \(r^{n,1}_t\) is the first-order natural real rate.
Second-order risk corrections will modify this object. The notation
\(r^{n,1}_t\) is used to keep the first-order natural rate conceptually separate
from the risk-adjusted natural rate.

\begin{remark}
Equations \eqref{eq:baseline_is}--\eqref{eq:baseline_taylor} are already
log-linear. Therefore, uncertainty effects do not arise by simply changing the
variance of \(a_t\) inside these equations. The risk terms must come from a
second-order expansion of the nonlinear household and firm conditions that
generate the reduced-form NK system.
\end{remark}

\section{TFP Level and TFP Uncertainty}

Let \(a_t\) be log TFP. The level innovation to TFP has stochastic volatility:
\begin{align}
    a_{t+1}
    &=
    \rho_a a_t
    +
    \exp(\sigma_{a,t})\varepsilon_{a,t+1},
    \qquad
    \varepsilon_{a,t+1}\sim N(0,1),
    \label{eq:tfp_level}
    \\
    \sigma_{a,t+1}
    &=
    (1-\rho_\sigma)\bar{\sigma}_a
    +
    \rho_\sigma \sigma_{a,t}
    +
    \sigma_\eta \eta_{\sigma,t+1},
    \qquad
    \eta_{\sigma,t+1}\sim N(0,1).
    \label{eq:tfp_vol}
\end{align}
The primitive first-moment shock is \(\varepsilon_{a,t+1}\). The primitive
uncertainty shock is \(\eta_{\sigma,t+1}\), which moves the conditional
volatility state \(\sigma_{a,t+1}\).

Conditional on information at date \(t\),
\begin{equation}
    \Var_t(a_{t+1})
    =
    \exp(2\sigma_{a,t}).
    \label{eq:conditional_tfp_variance}
\end{equation}
Because \(a_t\) is predetermined at date \(t\),
\begin{equation}
    \Var_t(\Delta a_{t+1})
    =
    \Var_t(a_{t+1})
    =
    \exp(2\sigma_{a,t}).
    \label{eq:conditional_tfp_growth_variance}
\end{equation}
Around the steady volatility level \(\bar{\sigma}_a\), the variance term can be
approximated by
\begin{equation}
    \exp(2\sigma_{a,t})-\exp(2\bar{\sigma}_a)
    \approx
    2\exp(2\bar{\sigma}_a)(\sigma_{a,t}-\bar{\sigma}_a).
    \label{eq:variance_linearization}
\end{equation}
Thus, even if the model is written in terms of variance wedges, the volatility
state can be used as the accounting state after local linearization.

Define the centered variance state
\begin{equation}
    q_t
    \equiv
    \exp(2\sigma_{a,t})-\exp(2\bar{\sigma}_a).
    \label{eq:variance_state}
\end{equation}
Under the single-uncertainty-source restriction, all uncertainty wedges should
be functions of \(q_t\), up to approximation error and endogenous amplification.

\section{Second-Order Expansion and Risk Terms}

\subsection{IS Equation}

The nonlinear household Euler equation can be written as
\begin{equation}
    1
    =
    \E_t
    \left[
        \beta
        \frac{U_c(C_{t+1})}{U_c(C_t)}
        \frac{R_t}{\Pi_{t+1}}
    \right].
    \label{eq:nonlinear_euler}
\end{equation}
Let
\begin{equation}
    m_{t+1}
    \equiv
    \log\beta
    +
    \log U_c(C_{t+1})
    -
    \log U_c(C_t)
    +
    i_t
    -
    \pi_{t+1}.
    \label{eq:log_sdf_return}
\end{equation}
Then Equation \eqref{eq:nonlinear_euler} is
\begin{equation}
    \E_t[\exp(m_{t+1})]=1.
\end{equation}
A second-order log-normal approximation gives
\begin{equation}
    0
    \approx
    \E_t m_{t+1}
    +
    \frac{1}{2}\Var_t(m_{t+1}).
    \label{eq:second_order_euler}
\end{equation}
With CRRA utility and local notation, this implies an intertemporal demand
condition of the form
\begin{equation}
    x_t
    =
    \E_t x_{t+1}
    -
    \frac{1}{\gamma}
    \left(
        i_t-\E_t\pi_{t+1}-r^{n,1}_t
    \right)
    +
    \omega^x_t,
    \label{eq:is_with_wedge}
\end{equation}
where the demand wedge collects second-order risk corrections:
\begin{equation}
    \omega^x_t
    \equiv
    -\frac{1}{2\gamma}
    \left[
        \Var_t(m_{t+1})
        -
        \Var^n_t(m^n_{t+1})
    \right]
    +
    \Delta r^n_{t,\mathrm{risk}}.
    \label{eq:demand_wedge_definition}
\end{equation}
Here \(m^n_{t+1}\) is the flexible-price counterpart and
\(\Delta r^n_{t,\mathrm{risk}}\) denotes the correction needed when the natural
real rate is defined as a risk-adjusted object rather than as the first-order
object \(r^{n,1}_t\).

The sign in Equation \eqref{eq:demand_wedge_definition} reflects precautionary
saving: for a given real interest rate path, higher conditional marginal-utility
risk lowers current demand relative to expected future demand.

\subsection{Phillips Curve}

The New Keynesian Phillips curve is generated by a nonlinear price-setting
condition. Write that condition generically as
\begin{equation}
    F^\pi
    \left(
        \pi_t,
        x_t,
        \E_t \pi_{t+1},
        \E_t x_{t+1},
        a_t,
        \Omega_t
    \right)
    =
    0,
    \label{eq:generic_price_setting}
\end{equation}
where \(\Omega_t\) is the conditional covariance matrix of future states. A
second-order expansion around the deterministic steady state yields the usual
linear NKPC plus a curvature correction:
\begin{equation}
    \pi_t
    =
    \beta \E_t \pi_{t+1}
    +
    \kappa x_t
    +
    \omega^\pi_t.
    \label{eq:nkpc_with_wedge}
\end{equation}
The markup or price-setting wedge can be represented as
\begin{equation}
    \omega^\pi_t
    =
    \Lambda_\pi
    \left[
        \frac{1}{2}
        \tr
        \left(
            H^\pi \Omega_t
        \right)
    \right],
    \label{eq:price_wedge_definition}
\end{equation}
where \(H^\pi\) is the Hessian of the underlying price-setting condition with
respect to future states, and \(\Lambda_\pi\) converts the curvature term into
inflation units.

Equation \eqref{eq:price_wedge_definition} is intentionally generic. Without
choosing Calvo versus Rotemberg pricing, the exact sign and magnitude of
\(\omega^\pi_t\) cannot be pinned down. The important accounting point is that
TFP uncertainty can appear in the NKPC only through the curvature of price
setting, real marginal cost, and expectations.

\subsection{Policy Rule}

If the monetary policy rule is linear and contains no policy uncertainty, a TFP
volatility shock does not create a primitive policy wedge:
\begin{equation}
    \omega^i_t = 0.
    \label{eq:no_policy_wedge}
\end{equation}
For empirical accounting, however, it is useful to allow the rule to be written
as
\begin{equation}
    i_t
    =
    \rho_i i_{t-1}
    +
    (1-\rho_i)
    \left(
        \phi_\pi \pi_t+\phi_x x_t
    \right)
    +
    \omega^i_t.
    \label{eq:taylor_with_wedge}
\end{equation}
Under the single-source TFP uncertainty model, \(\omega^i_t\) should be zero
unless it is interpreted as a residual that absorbs rule misspecification,
policy smoothing not captured by the simple rule, or measurement error.

\section{Wedge Representation}

Collect the three accounting wedges as
\begin{equation}
    \bs{\omega}_t
    =
    \begin{bmatrix}
        \omega^x_t \\
        \omega^\pi_t \\
        \omega^i_t
    \end{bmatrix}.
    \label{eq:wedge_vector}
\end{equation}
The wedge-augmented NK system is
\begin{align}
    x_t
    &=
    \E_t x_{t+1}
    -
    \frac{1}{\gamma}
    \left(
        i_t-\E_t\pi_{t+1}-r^{n,1}_t
    \right)
    +
    \omega^x_t,
    \label{eq:wedge_system_is}
    \\
    \pi_t
    &=
    \beta \E_t \pi_{t+1}
    +
    \kappa x_t
    +
    \omega^\pi_t,
    \label{eq:wedge_system_nkpc}
    \\
    i_t
    &=
    \rho_i i_{t-1}
    +
    (1-\rho_i)
    \left(
        \phi_\pi \pi_t+\phi_x x_t
    \right)
    +
    \omega^i_t.
    \label{eq:wedge_system_policy}
\end{align}

Because TFP uncertainty is the only primitive source of volatility, the first
restricted accounting representation is
\begin{equation}
    \bs{\omega}_t
    =
    \bs{\lambda} q_t
    +
    \bs{u}_t,
    \qquad
    \bs{\lambda}
    =
    \begin{bmatrix}
        \lambda_x \\
        \lambda_\pi \\
        \lambda_i
    \end{bmatrix},
    \label{eq:rank_one_wedge_restriction}
\end{equation}
where \(q_t\) is defined in Equation \eqref{eq:variance_state} and
\(\bs{u}_t\) is approximation error or residual model misspecification.

The baseline theoretical restrictions are:
\begin{align}
    \lambda_x &< 0
    && \text{if precautionary saving dominates,}
    \label{eq:lambda_x_prediction}
    \\
    \lambda_\pi &\ \text{ambiguous}
    && \text{until the price-setting microfoundation is fixed,}
    \label{eq:lambda_pi_prediction}
    \\
    \lambda_i &= 0
    && \text{under a linear Taylor rule with no policy uncertainty.}
    \label{eq:lambda_i_prediction}
\end{align}

This representation is useful because it provides a sharp diagnostic. If a
single TFP uncertainty shock is the correct primitive force, then the three
wedges should be highly collinear and mostly driven by \(q_t\). If the inferred
wedges are not well described by Equation \eqref{eq:rank_one_wedge_restriction},
then either the simplified NK model is missing propagation mechanisms or more
than one primitive uncertainty source is needed.

\begin{table}[htbp]
\centering
\caption{Interpretation of NK Uncertainty Wedges}
\label{tab:wedge_interpretation}
\begin{tabular}{lll}
\toprule
Wedge & Equation & Interpretation \\
\midrule
\(\omega^x_t\) & IS curve & Precautionary saving, risk-adjusted real rate, demand risk \\
\(\omega^\pi_t\) & NKPC & Markup risk, price-setting curvature, marginal-cost risk \\
\(\omega^i_t\) & Taylor rule & Policy residual; zero in the baseline theory \\
\bottomrule
\end{tabular}
\end{table}

\section{BCA-Style Accounting Logic}

The accounting exercise follows the spirit of Business Cycle Accounting. The
goal is not to claim that a wedge is itself a deep friction. The goal is to
measure which equilibrium condition must be distorted in order for the simple
model to reproduce the observed dynamics.

\subsection{Step 1: Infer Wedges}

Given observed or model-implied series for
\(\{x_t,\pi_t,i_t,r^{n,1}_t\}\) and expectations
\(\{\E_t x_{t+1},\E_t \pi_{t+1}\}\), infer the wedges as residuals:
\begin{align}
    \widehat{\omega}^x_t
    &=
    x_t
    -
    \E_t x_{t+1}
    +
    \frac{1}{\gamma}
    \left(
        i_t-\E_t\pi_{t+1}-r^{n,1}_t
    \right),
    \label{eq:infer_x_wedge}
    \\
    \widehat{\omega}^\pi_t
    &=
    \pi_t
    -
    \beta \E_t\pi_{t+1}
    -
    \kappa x_t,
    \label{eq:infer_pi_wedge}
    \\
    \widehat{\omega}^i_t
    &=
    i_t
    -
    \rho_i i_{t-1}
    -
    (1-\rho_i)(\phi_\pi\pi_t+\phi_x x_t).
    \label{eq:infer_i_wedge}
\end{align}

\subsection{Step 2: Test the Single-Uncertainty-Source Restriction}

Estimate or inspect the relation
\begin{equation}
    \widehat{\bs{\omega}}_t
    =
    \bs{\lambda} q_t
    +
    \bs{u}_t.
    \label{eq:estimate_wedge_loadings}
\end{equation}
The single-source TFP uncertainty hypothesis implies that most systematic
variation in the wedge vector should be explained by the scalar variance state
\(q_t\). In this minimal model, a large and systematic \(\widehat{\omega}^i_t\)
would be evidence against the clean baseline, because monetary policy is not a
primitive uncertainty source.

\subsection{Step 3: Counterfactual Wedge Experiments}

Once the wedges are inferred, run counterfactual paths by opening one wedge at a
time:
\begin{enumerate}[label=(\alph*)]
    \item Demand-risk counterfactual: set
    \(\omega^x_t=\widehat{\omega}^x_t\) and
    \(\omega^\pi_t=\omega^i_t=0\).
    \item Markup-risk counterfactual: set
    \(\omega^\pi_t=\widehat{\omega}^\pi_t\) and
    \(\omega^x_t=\omega^i_t=0\).
    \item Policy-residual counterfactual: set
    \(\omega^i_t=\widehat{\omega}^i_t\) and
    \(\omega^x_t=\omega^\pi_t=0\).
\end{enumerate}
The contribution of each wedge is evaluated by how closely the counterfactual
paths reproduce the observed paths of \(x_t\), \(\pi_t\), and \(i_t\).

\subsection{Connection to the Later SOE Nonlinear BCA}

The later small-open-economy framework will contain more wedges: efficiency,
labor, investment, trade, and UIP wedges. This note is a stripped-down
laboratory. Its value is to establish the logic:
\[
    \text{second-moment primitive}
    \rightarrow
    \text{equation-specific risk wedge}
    \rightarrow
    \text{counterfactual accounting}.
\]
Once this logic is clear, the same structure can be expanded to multiple
external uncertainty factors and multiple small-open-economy wedges.

\section{Open Questions}

This version leaves several questions intentionally open.

\begin{enumerate}[label=\textbf{Q\arabic*.}]
    \item \textbf{Second order or third order?}
    The second-order expansion identifies risk correction terms in equilibrium
    equations. However, a fully solved stochastic-volatility DSGE model usually
    requires third-order perturbation for volatility shocks to affect decision
    rules as independent state variables. The next step is to decide whether the
    paper needs a full third-order solution or only an accounting
    representation.

    \item \textbf{Where should wedges be defined?}
    This note defines wedges in the reduced-form IS curve, NKPC, and Taylor
    rule. A deeper BCA exercise could instead define wedges at the level of
    household Euler equations, firm pricing FOCs, and resource constraints.

    \item \textbf{Which pricing model pins down \(\lambda_\pi\)?}
    The sign of the price-setting wedge depends on the nonlinear pricing
    environment. Calvo and Rotemberg pricing can imply different uncertainty
    transmission mechanisms, so the markup wedge should not be interpreted
    before choosing the pricing microfoundation.

    \item \textbf{Should natural-rate risk be absorbed into \(\omega^x_t\)?}
    One can define a risk-adjusted natural rate and keep the IS wedge smaller,
    or define the natural rate at first order and place all risk corrections in
    \(\omega^x_t\). The latter is cleaner for accounting; the former may be
    cleaner for structural interpretation.

    \item \textbf{How will this connect to the empirical uncertainty factors?}
    This note has one volatility state \(q_t\). The larger research project has
    empirical macro and financial uncertainty factors. The future SOE version
    will need a factor-loading structure that maps those empirical factors into
    multiple structural wedge volatilities.
\end{enumerate}


\appendix
\section{Generic Algebra for a Second-Order Accounting Wedge}

This appendix records the generic algebra behind the wedge construction. Let an
equilibrium condition be written as
\begin{equation}
    \E_t
    \left[
        F(z_t,z_{t+1})
    \right]
    =
    0,
    \label{eq:generic_equilibrium_condition}
\end{equation}
where \(z_t\) stacks the relevant endogenous and exogenous variables. Let
\(\bar{z}\) be the deterministic steady state and define
\(\tilde{z}_t=z_t-\bar{z}\). A second-order expansion gives
\begin{equation}
    0
    \approx
    F(\bar{z},\bar{z})
    +
    F_0 \tilde{z}_t
    +
    F_1 \E_t \tilde{z}_{t+1}
    +
    \frac{1}{2}
    \E_t
    \left[
        \tilde{z}_{t+1}'
        H_1
        \tilde{z}_{t+1}
    \right]
    +
    \text{terms involving } \tilde{z}_t^2.
    \label{eq:generic_second_order_expansion}
\end{equation}
At the deterministic steady state, \(F(\bar{z},\bar{z})=0\). The conditional
quadratic term can be decomposed as
\begin{equation}
    \E_t
    \left[
        \tilde{z}_{t+1}'
        H_1
        \tilde{z}_{t+1}
    \right]
    =
    (\E_t\tilde{z}_{t+1})'
    H_1
    (\E_t\tilde{z}_{t+1})
    +
    \tr(H_1\Omega_t),
    \label{eq:mean_variance_decomposition}
\end{equation}
where
\begin{equation}
    \Omega_t
    =
    \Var_t(\tilde{z}_{t+1}).
    \label{eq:omega_covariance}
\end{equation}

For an accounting exercise focused on uncertainty, collect all terms involving
\(\Omega_t\) into a wedge:
\begin{equation}
    \omega^F_t
    \equiv
    -A_F^{-1}
    \left[
        \frac{1}{2}\tr(H_1\Omega_t)
    \right],
    \label{eq:generic_wedge}
\end{equation}
where \(A_F\) is the coefficient on the endogenous variable being isolated. This
is the general form behind Equations \eqref{eq:demand_wedge_definition} and
\eqref{eq:price_wedge_definition}.

In the TFP-volatility-only case, the relevant element of \(\Omega_t\) is
\(\exp(2\sigma_{a,t})\). If the rest of the covariance matrix is fixed or
endogenously proportional to this variance state, then the wedge has the local
representation
\begin{equation}
    \omega^F_t
    \approx
    \lambda_F
    \left[
        \exp(2\sigma_{a,t})-\exp(2\bar{\sigma}_a)
    \right].
    \label{eq:generic_rank_one_wedge}
\end{equation}
This is the algebraic source of the rank-one restriction in Equation
\eqref{eq:rank_one_wedge_restriction}.


\end{document}
